\subsection{The full congruence quotient together with abelianization}

\begin{statement}{A.56}{Proved}{Exact abelian image of the full congruence kernel}
Let $H$ be the full geometric Legendre image in A.17, and put
\[
 C=\ker(\Pi\longrightarrow H),\qquad
 T=\Pi^{\rm ab}=\widehat\Z(1)^2.
\]
All derived subgroups in this statement are closed derived
subgroups.  There is an isomorphism of finite $G_\Q$-modules
\[
 H^{\rm ab}\simeq (\Z/8(1))^2\oplus\Z/3(1),
\tag{A.56.1}
\]
under which the map induced by the two Kummer coordinates is
\[
 T\longrightarrow H^{\rm ab},\qquad
          (a,b)\longmapsto(a\bmod8, b\bmod8, a+b\bmod3).
\tag{A.56.2}
\]
In particular, the image of $C$ in $T$ is exactly the open subgroup
\[
 K=\{(a,b):a\equiv b\equiv0\pmod8,\quad a+b\equiv0\pmod3\}.
\tag{A.56.3}
\]

Set $E_H=E_U/C$, $E_{H^{\rm ab}}=E_H/[H,H]$, and
$E_1=E_U/\Gamma_2\Pi$.  The natural map is an isomorphism of
profinite extensions,
\[
 E_U/(C\cap\Gamma_2\Pi)
       \xrightarrow{\ \sim\ }
       E_H\mathop{\times}_{E_{H^{\rm ab}}}E_1.
\tag{A.56.4}
\]
Fix a rational point $b_0\in U(\Q)$.  Given completed Kummer
classes $(A,B)$, a section of (A.56.4) with these abelian classes
and with its $E_H$ image equal to the point section at $b_0$
exists exactly when
\[
 \begin{split}
 A&\equiv b_0\pmod{(\widehat{\Q^\times})^8},\\
 B&\equiv1-b_0\pmod{(\widehat{\Q^\times})^8},\\
 AB&\equiv b_0(1-b_0)\pmod{(\widehat{\Q^\times})^3}.
 \end{split}
\tag{A.56.5}
\]
These expressions mean equality under the canonical projections
to $\Q^\times/(\Q^\times)^8$ and
$\Q^\times/(\Q^\times)^3$, respectively.

The ordinary infinite-support ghosts of A.44 satisfy (A.56.5)
for $b_0=3$.  They therefore define global sections of this
quotient with the entire congruence representation of the honest
point $3$, including its determinant and its integral pairings.
At every fresh support prime in the construction, these quotient
sections fail the local Legendre comparison with their prescribed
ordinary point.  They are therefore not images of Selmer sections
of $E_U$.  The exact full-section recognition question in A7 is
not decided by this quotient construction.
\end{statement}
\begin{proof}
We compute the closed derived subgroups of the factors of $H$.
At the prime two write
\[
 H_2=\left\{\begin{pmatrix}a&b\\c&d\end{pmatrix}
   \in\operatorname{SL}_2(\Z_2):
       b,c\in2\Z_2,\ a,d\equiv1\pmod4\right\},
 \qquad x=U(2),\quad y=L(-2).
\tag{A.56.6}
\]
These two matrices topologically generate $H_2$ by A.17.
Let $D_2=[H_2,H_2]$ and
$K_n=\ker(\operatorname{SL}_2(\Z_2)\to
\operatorname{SL}_2(\Z/2^n))$.  The exact integer commutator is
\[
                c=[x,y]=\begin{pmatrix}13&8\\8&5\end{pmatrix}.
\tag{A.56.7}
\]
Modulo 16 it is central, has square $9I$, and has order four.
The group $H_2$ modulo 16 has order 256: its reduction modulo
four has four elements, and each of the next two reduction
maps has kernel of order eight.  Therefore its derived subgroup
is precisely $\langle c\rangle$ of order four, and its
abelianization has order 64.  Both $x,y$ have order eight
modulo 16.  They generate that abelianization, so the map
\[
         (\Z/8)^2\longrightarrow
            (H_2\bmod16)/\langle c\rangle,
                         (a,b)\longmapsto x^a y^b
\tag{A.56.8}
\]
is an isomorphism.

We next prove that no further abelian quotient appears at
higher powers of two.  The following three elements of $D_2$
are in $K_4$, with the indicated leading terms modulo 32:
\[
 \begin{array}{c|c}
 \text{element}&(g-I)/16\pmod2\\ \hline
 c^4&I\\
 {[x,c]}&I+e_{12}\\
 {[y,c]}&I+e_{21}.
 \end{array}
\tag{A.56.9}
\]
For example their exact integer matrices are
\[
 \begin{pmatrix}75025&46368\\46368&28657\end{pmatrix},\quad
 \begin{pmatrix}465&-752\\128&-207\end{pmatrix},\quad
 \begin{pmatrix}81&-128\\-112&177\end{pmatrix}.
\]
The three leading terms span
$\mathfrak{sl}_2(\mathbf F_2)$, which is
$K_4/K_5$ under $g\mapsto(g-I)/16$.  If $n\ge4$, then
\[
 (I+2^nA)^2\equiv I+2^{n+1}A\pmod{2^{n+2}}.
\tag{A.56.10}
\]
Repeated squares of the three elements consequently show that
$D_2\cap K_n$ surjects onto $K_n/K_{n+1}$ for every $n\ge4$.
For any element of $K_4$, choose successively products of
these elements to match it modulo $2^5,2^6,\ldots$.
The corrections tend to the identity and the products converge
in $D_2$, which is closed.  Their limit is the prescribed
element.  Hence $K_4\subset D_2$.
Together with (A.56.7)--(A.56.8) this proves
\[
 D_2=\{g\in H_2:g\bmod16\in\langle c\rangle\},
                         \qquad H_2^{\rm ab}=(\Z/8)^2.
\tag{A.56.11}
\]
This proof controls the full inverse limit rather than only
the first finite reductions.

At three put $G_3=\operatorname{SL}_2(\Z_3)$.
The commutator formula
\[
 [\operatorname{diag}(2,2^{-1}),U(a)]=U(3a)
\tag{A.56.12}
\]
and the corresponding lower-triangular formula show that
$U(3\Z_3)$ and $L(3\Z_3)$ lie in $[G_3,G_3]$.
For $d\in1+3\Z_3$, the identity
\[
 \operatorname{diag}(d,d^{-1})=
 U(1)L(d-1)U(-1/d)L(-d(d-1))
\tag{A.56.13}
\]
also has trivial image in the abelianization: the two lower
factors already do, and the two upper factors have product
$U(1-1/d)$ there, with $1-1/d\in3\Z_3$.
Gauss decomposition now puts the entire reduction kernel
modulo three in the closed derived subgroup.

The derived subgroup of $\operatorname{SL}_2(\mathbf F_3)$
is the quaternion group of order eight.  Explicitly, the set
consisting of $\pm I$ and the six trace-zero matrices is a
normal quaternion subgroup, with quotient of order three.
The commutator of $U(1)$ and $L(1)$ has order four, and its
conjugates by these generators generate that quaternion
subgroup.  Thus
\[
                G_3^{\rm ab}=\Z/3.
\tag{A.56.14}
\]
Moreover $U(2)$ and $L(-2)$ have the same nonzero image in
this quotient, since
\[
 S U(2)S^{-1}=L(-2),\qquad
                      S=\begin{pmatrix}0&-1\\1&0\end{pmatrix}.
\tag{A.56.15}
\]
The image is nonzero because either unipotent has order three
modulo three and is outside the quaternion subgroup.

For every $\ell\ge5$ there is a unit $u\in\Z_\ell^\times$
with $u^2-1\in\Z_\ell^\times$.  The commutator of
$\operatorname{diag}(u,u^{-1})$ with $U(a)$ is
$U((u^2-1)a)$, and similarly for lower matrices.  Thus the
closed derived subgroup contains all elementary matrices,
which generate $\operatorname{SL}_2(\Z_\ell)$ by row
reduction.  All these factors are perfect.
The closed derived subgroup of a product of these profinite
groups is the product of their closed derived subgroups:
the inclusion in one direction follows by projection, and
in the other direction one first realizes the elements in
finitely many factors and then takes limits.  This proves
the group structure in (A.56.1).

We specify its character map, including the sign at three.
Take positive cusp meridians $\gamma_0,\gamma_1$ whose
Kummer degrees for $(t,1-t)$ are $(1,0),(0,1)$.
In the labeled convention compatible with A.16 the ordered
pair of Picard--Lefschetz matrices can be chosen as
\[
                \rho(\gamma_0)=L(-2),\qquad
                \rho(\gamma_1)=U(2).
\tag{A.56.16}
\]
Equivalently one may simultaneously conjugate by $S$ and
interchange the two matrix names.  Changes of paths to the
cusps conjugate the individual meridians and do not change
their abelian images.  At two use (A.56.8), with its two
coordinates ordered as $(\gamma_0,\gamma_1)$, to obtain
$(a\bmod8,b\bmod8)$.  At three (A.56.15) gives $a+b$,
and not $a-b$.  The same plus sign is independently fixed by
the actual cubic quotient $v^3=4t(1-t)$ in A.39: both finite
cusps have valuation one in its defining unit.
Finally the map $T\to H^{\rm ab}$ is $G_\Q$-equivariant.
The arithmetic action on $T$ is cyclotomic in both
coordinates, so its quotient action on the displayed
factors is precisely that in (A.56.1).  This establishes
the full module map (A.56.2).

We now prove the assertion about the full congruence kernel.
The map from the compact group $\Gamma_2\Pi$ to $[H,H]$
is onto.  Its image is closed, is contained in $[H,H]$,
and contains every commutator of $H$ by lifting the two
entries to $\Pi$; it therefore contains their closed
generated subgroup.  If an element of $\Pi$ has abelian
image $t\in T$ mapping to zero in $H^{\rm ab}$, its image
in $H$ lies in $[H,H]$.  Multiply it by a suitable element
of $\Gamma_2\Pi$ to cancel this image.  The result belongs
to $C$ and still has abelian image $t$.  Conversely every
element of $C$ has zero image in $H^{\rm ab}$.
The image of $C$ in $T$ is therefore exactly (A.56.3).

The same lifting argument proves that
\[
                \Pi\longrightarrow H\times_{H^{\rm ab}}T
\tag{A.56.17}
\]
is onto: lift the $H$ coordinate first, and then correct
its $T$ coordinate by an element of $C$.  The kernel is
$C\cap\Gamma_2\Pi$.  On each arithmetic fibre over
$G_\Q$ apply the same argument, first choosing one lift
in $E_U$ of its specified $E_H$ coordinate.
This proves surjectivity and the claimed kernel in
(A.56.4).  All maps are continuous and all groups compact
Hausdorff, so the resulting bijection is a topological
isomorphism.  In particular the passage from finite $Q$
in A.51 to the full $H$ has been proved, rather than
inferred from separate finite-level splittings.

Use the point $b_0$ to split all these arithmetic
extensions.  Relative to its abelian splitting, $(A,B)$
gives a continuous cocycle $z:G_\Q\to T$ representing
\[
               (A/\kappa(b_0),\ B/\kappa(1-b_0)).
\tag{A.56.18}
\]
Its image in $H^1(G_\Q,H^{\rm ab})$ is zero exactly when
(A.56.5) holds, by (A.56.1)--(A.56.2) and the Kummer
isomorphisms for $\mu_8$ and $\mu_3$.
Vanishing of this cohomology class means that the image
cocycle is a coboundary, say $\sigma v-v$ for
$v\in H^{\rm ab}$.  Lift $v$ through the surjective
map $T\to H^{\rm ab}$ and conjugate the abelian section
by this lift, with the sign that cancels that coboundary.
Its image in $E_{H^{\rm ab}}$ is now literally the point
section at $b_0$.  Pairing it with the point section in
$E_H$ and applying (A.56.4) constructs the required
continuous quotient section.  Conversely a section
with the prescribed $E_H$ component gives zero image
of (A.56.18) in this cohomology group, and hence
(A.56.5).  The conjugation changes no Kummer cohomology
class, so this proves the assertion for the specified
$(A,B)$, including its necessity.

In A.44, $A$ and $B$ are obtained from $3$ and $-2$ by
multiplication by prime powers with exponents
$n_j=(j+4)!$.  All these exponents are divisible by
$24$, the exponent of (A.56.1).  Their limits therefore
satisfy (A.56.5) for $b_0=3$.
The construction above gives their quotient section
with the complete $E_H$ image of that point.

We finish by checking exactly why its ordinary local
data do not have that image.  At a fresh prime $q_j$
in the support of $A$, A.44 gives
\[
                v_{q_j}(x_{q_j})=n_j>0,
                    v_{q_j}(1-x_{q_j})=0.
\tag{A.56.19}
\]
All these primes are distinct from two and three.
The Legendre discriminant and invariant are
\[
       \Delta(t)=16t^2(1-t)^2,\qquad
                    c_4(t)=16(t^2-t+1).
\]
Thus $E_{x_{q_j}}$ has multiplicative reduction at
$q_j$ with discriminant valuation $2n_j>0$.
For every coefficient prime $\ell\ne q_j$, after an
unramified quadratic splitting extension its inertia
action on $T_\ell$ is a unipotent whose off-diagonal
tame character has coefficient $2n_j$.  This is
nonzero in $\Z_\ell$ even when it is highly divisible
by $\ell$.
At a fresh prime $r_j$ in the support of $B$ the
same calculation applies with $x$ and $1-x$
interchanged.

By contrast $E_3$ has discriminant $576$ and good
reduction at every fresh support prime.  Its
$\ell$-adic inertia action there is trivial for
$\ell$ different from that prime.  The quotient
section in (A.56.4) has exactly this representation
because its $E_H$ image is the point section at $3$.
Nontrivial unipotent inertia cannot be conjugated
to trivial inertia.  This proves the claimed local
failure in the full congruence quotient.

There is no alternative ordinary local point with
the same abelian coordinates: the Kummer map
$\Q_q^\times\to\widehat{\Q_q^\times}$ is injective,
as follows from its decomposition into the valuation
factor, finite roots of unity, and principal units.
It fixes that point as $x_q$.  Nor can a cuspidal
section have these coordinates.  At a fresh
$q_j$, cusp zero would require
$\kappa(1-x_{q_j})=1$, impossible by injectivity
and $x_{q_j}\ne0$; cusp one and infinity are
excluded by the distinct valuations $(n_j,0)$.
The analogous argument at $r_j$ exchanges zero
and one.  Thus the quotient sections are not
local point or cusp images there, and in particular
do not lift to Selmer sections of $E_U$.
This last verification identifies the additional
local condition required by the full A7 problem.
\end{proof}

\noindent\emph{Exact cyclic and Fermat comparison.}
For integers $a,b$ and $n\ge1$, the Kummer character of the
finite etale cover
\[
                         z^n=t^a(1-t)^b
\]
factors through $E_H$ if and only if
\[
                   n\mid24a,\qquad n\mid24b,
                                  \qquad n\mid8(a-b).
\tag{A.56.20}
\]
Indeed its restriction to $T$ is $(u,v)\mapsto au+bv\pmod n$,
and it factors exactly when it kills the image $K$ of $C$.
The module $K$ in (A.56.3) is generated over $\widehat\Z$ by
$(24,0),(0,24),(8,-8)$: write an element as $(8u,8v)$;
then $u+v=3w$, and
$(8u,8v)=w(0,24)+u(8,-8)$.
Evaluation on these generators proves (A.56.20).
The full Fermat cover $x^n=t$, $y^n=1-t$ therefore factors
through the congruence quotient exactly when $n\mid8$,
by applying (A.56.20) to its two coordinate characters.
This supplies the precise map and the exact range needed
when comparing Fermat covers with congruence covers in L15.

\noindent\emph{The common quotient with a geometric pro-$\ell$ group.}
Let $N_\ell=\ker(\Pi\to\Pi^{(\ell)})$, with
$\Pi^{(\ell)}$ the maximal pro-$\ell$ quotient.
The maximal pro-$\ell$ quotients of $H$ are exactly
\[
 H^{(\ell)}=
 \begin{cases}
 H_2,&\ell=2,\\
 \Z/3,&\ell=3,\\
 1,&\ell\ge5.
 \end{cases}
\tag{A.56.21}
\]
To verify this, every factor $\operatorname{SL}_2(\Z_p)$
for $p\ge5$ is perfect, so it has no nontrivial finite
$\ell$-group quotient: a nontrivial finite $\ell$-group
always has nontrivial abelianization.  The factor $H_2$
is pro-two.  The pro-three quotient of
$\operatorname{SL}_2(\Z_3)$ has abelianization $\Z/3$
by (A.56.14), and is therefore generated by one element.
Indeed, if the closed cyclic subgroup generated by a lift
of its abelian generator were proper, a finite three-group
quotient would place it in a maximal subgroup of index
three, contradicting that choice of generator.
This quotient is consequently cyclic and equals $\Z/3$.
The same maximal-subgroup argument excludes every pro-two
quotient of this factor, since its abelianization has odd
order.  Its finite quotients have orders divisible only
by two and three, which excludes all other coefficient
primes.  These arguments and continuity on the product
prove (A.56.21).

The image of $N_\ell$ in $H$ is the kernel of
$H\to H^{(\ell)}$.  One inclusion follows by composing
with that pro-$\ell$ quotient, and the other follows
because quotienting $H$ by the image of $N_\ell$ gives
a pro-$\ell$ quotient of $\Pi$ and hence of $H$.
Apply the same compact fibre-product proof as in
(A.56.17) to the two closed normal subgroups
$C,N_\ell$.  For $\ell\ge5$ it gives the exact
isomorphism of arithmetic extensions
\[
 E_U/(C\cap N_\ell)
       \xrightarrow{\ \sim\ }
                  E_H\times_{G_\Q}(E_U/N_\ell).
\tag{A.56.22}
\]
Thus any section of the geometrically pro-$\ell$
quotient can be paired globally with a prescribed
point section in the full congruence quotient.
For $\ell=3$ the fibre product is instead over the
arithmetic cubic quotient of A.39, whose geometric
character is $a+b\bmod3$; for $\ell=2$ it is over
the quotient with geometric group $H_2$.
These are actual quotient maps.  As in the preceding
construction, the assertion about the global fibre
product does not prescribe the joint local point
images or a section of the full group $E_U$.

\noindent\emph{Retained verification.}
\texttt{code/legendre\_congruence\_abelian.py} verifies the
integer matrices and leading terms in (A.56.7)--(A.56.9), the
unique $x^a y^b c^k$ forms modulo 16, all products in their
abelian quotient, and the mod-three and mod-nine calculations.
Its complete output is retained in
\texttt{code/legendre\_congruence\_abelian\_output.txt}.
The proof above supplies the passage from these finite
certificates to the closed derived subgroups and the full
arithmetic fibre product.
